\setcounter{section}{29}

  \zwischenueberschrift{Soal latihan}

\inputaufgabe
{}
{

Misalkan $M$ suatu
\definitionsverweis {manifold diferensiabel}{}{}
sedemikian sehingga
\definitionsverweis {bundel tangen}{}{}
$TM$ bersifat trivial. Melalui suatu trivialisasi diferensiabel

\mavergleichskette
{\vergleichskette
{ TM
}
{ \cong }{ M \times \R^n
}
{  }{
}
{    }{
}
{   }{
}
}
{}{}{}
lengkapi $M$
\zusatzklammer {dengan menggunakan
\definitionsverweis {hasil kali dalam standar}{}{}
pada $\R^n$} {} {}
dengan suatu
\definitionsverweis {struktur Riemann}{}{}
\zusatzklammer {lihat
Soal 16.2} {} {.}
Tunjukkan bahwa
\definitionsverweis {operator kelengkungan}{}{}
pada $M$ bersifat trivial.

}
{} {}

\inputaufgabegibtloesung
{}
{

Dalam situasi dua dimensi pada
Lemma 29.1,
hitung simbol-simbol Christoffel yang relevan untuk
\definitionsverweis {operator kelengkungan}{}{.}

}
{} {}

\inputaufgabe
{}
{

Kita tinjau
\definitionsverweis {torus}{}{}
$S^1 \times S^1$. Pertama, lengkapi torus tersebut dengan struktur yang diberikan oleh struktur produk
\zusatzklammer {bersama penyematan

\mavergleichskette
{\vergleichskette
{ S^1
}
{ \subseteq }{ \R^2
}
{  }{
}
{    }{
}
{   }{
}
}
{}{}{,}
lihat
Soal 16.3
dan
Soal 16.4} {} {}
dan
\definitionsverweis {struktur Riemann}{}{,}
kemudian bandingkan dengan realisasi tertanam berikut.

\mavergleichskettedisp
{\vergleichskette
{ T
}
{ =} { \left\{ (x,y,z) \in \R^3  \mid   \left(  \sqrt{x^2+y^2} -R  \right)^2 +z^2  = r^2         \right\}
}
{ \subseteq} { \R^3
}
{ } {
}
{ } {
}
}
{}{}{}
dengan struktur Riemann terinduksi. Bandingkan
\definitionsverweis {operator-operator kelengkungan}{}{}
untuk kedua struktur tersebut.

}
{} {}

  \zwischenueberschrift{Soal untuk dikumpulkan}
